% \item---------------------------------------------------------
% Indentação
% ----------------------------------------------------------
\chapter[Fundamentação Teórica]{Fundamentação Teórica}

A Engenharia de Gestão (\textit{Management Engineering} -- ManagEng) caracteriza-se como a disciplina voltada à construção de modelos matemáticos de sistemas reais e à sua subsequente análise para subsidiar tomadas de decisão gerenciais fundamentadas \cite[p.~4]{kolker2012}. Diferentemente de abordagens puramente administrativas, essa área aplica princípios de engenharia e metodologias analíticas para compreender o impacto de novas tecnologias nas organizações e na sociedade, permitindo projetar sistemas de gestão integrados e processos de inovação que consideram a perspectiva do usuário, as mudanças ambientais e as interconexões de mercado \cite{duimering2013, elia2021}. Do ponto de vista profissional, o engenheiro de gestão atua frequentemente como um consultor interno de melhorias, assumindo responsabilidades que demandam o cultivo de uma visão sistêmica e holística das operações organizacionais \cite{larson2013}.

Na literatura científica e na prática profissional, observa-se uma recorrente ambiguidade terminológica entre Engenharia de Gestão (\textit{Management Engineering}) e Gestão da Engenharia (\textit{Engineering Management} -- EngManag), bem como uma frequente sobreposição com a Engenharia Industrial (\textit{Industrial Engineering}), especialmente no setor da saúde \cite{lannes2001, kotnour2005}. Embora tanto a Engenharia de Gestão quanto a Gestão da Engenharia atuem como pontes entre os campos da engenharia e da administração, ambas partem de perspectivas distintas \cite{elia2021}.

A Gestão da Engenharia é definida como o conjunto de conhecimentos, habilidades e responsabilidades gerenciais gerais que engenheiros adquirem à medida que progridem em suas carreiras e assumem cargos de liderança \cite{asem2019, shah2015}. O foco dessa vertente reside na alavancagem da experiência técnica de engenheiros para supervisionar pessoas, projetos e recursos de forma eficaz em ambientes tecnologicamente complexos \cite{mohn2022}. Em contrapartida, a Engenharia de Gestão oferece uma perspectiva de engenharia aplicada diretamente à própria função gerencial, concentrando-se na modelagem quantitativa, otimização, tomada de decisão, gestão de projetos, gestão de riscos e análise de sistemas \cite{kolker2012, duimering2013}. Enquanto a Gestão da Engenharia possui uma estrutura bibliométrica amplamente consolidada e mapeada na literatura acadêmica \cite{philbin2024, cardenas2023}, a Engenharia de Gestão ainda carece de reconhecimento formal como ocupação distinta nas principais taxonomias internacionais de trabalho, tais como a \textit{International Standard Classification of Occupations} (ISCO-08), a \textit{Occupational Information Network} (O*NET) e a \textit{European Skills, Competences, Qualifications and Occupations} (ESCO) \cite{ilo2012, onet, esco}.

Historicamente, a produção científica sobre Engenharia de Gestão remonta ao final da década de 1920, apresentando longos hiatos de publicação até meados do século XX, período em que estudos pioneiros exploraram a aplicação de princípios de engenharia e gestão na administração hospitalar e na enfermagem \cite{gilbreth1950, bartscht2010}. Nas décadas seguintes, a disciplina consolidou sua base teórica voltada ao planejamento, à modelagem interna de competências e à gestão de sistemas sociotécnicos \cite{bonjour2002, lal2013}, expandindo-se mais recentemente para frentes como gestão ambiental, economia verde, gestão de projetos e educação em engenharia \cite{liu2020, biancardi2023}.

Apesar de sua relevância prática e evolução histórica, a consolidação da Engenharia de Gestão como um campo científico autônomo e claramente delimitado em relação à Gestão da Engenharia representa um esforço analítico recente na literatura. Nesse sentido, destaca-se o mapeamento bibliométrico conduzido por \citeonline{karakiozis2026}, que analisou a evolução das estruturas sociais e conceituais da área entre 1928 e 2025, evidenciando as divergências conceituais entre países e reforçando a necessidade de distinguir formalmente a identidade acadêmica e profissional da Engenharia de Gestão.